\chapter{Needs and Strategies}

Feelings are rapid signals that tell us whether our needs are being met;
they serve as the bridge to understanding our inner world.
However, a feeling is merely a signal.
What truly drives this signal is the \textbf{need} behind it.
If feelings are the color and brightness of a light, then needs are the
energy source generating that light.
This chapter explores the essence of needs, the variety of universal human
needs, and how we can authentically identify and express our deepest needs
so that our communication genuinely touches the heart.

\section{Needs}

\subsection{Definition}

\subsubsection{The Source of Feelings}

Needs are the internal driving forces that we all share, which generate our
feelings.
Stepping away from dualistic thinking is also known as ``needs-based thinking.''
Feelings are not directly caused by the actions of others, but by whether
our own needs are being satisfied.
This is one of the core insights of Nonviolent Communication:
\textbf{No one can make you feel a certain way; your feelings come only
from your own needs} \citep{rosenberg2003}.
We take responsibility for our own needs, rather than pushing the
responsibility onto others or forcing ourselves to meet others' needs.
Let us return to the late-night homecoming scenario to
deeply understand the relationship between feelings and needs:

\paragraph{Case Analysis: The Husband Returning Late (Needs Version)}

The husband opens the door and comes home at 11 PM.
The wife has been waiting for a long time, and complex emotions are surging
within her.
She feels lonely and overwhelmed.
Is the husband's late return the ``cause'' of her loneliness?

On the surface, it seems so.
But if we ask a deeper question: \textbf{``She is lonely, because what need
is not being met?''}
The answer becomes clear: She is lonely because she deeply \textbf{needs
companionship, connection, and support} \citep{bowlby1969}.

This distinction is crucial:
\begin{itemize}
    \item \textbf{Judgmental Thinking (Alienating Communication):}
    ``You make me feel lonely, it's your fault, it's all because you didn't
    fulfill your duties as a husband.''
    (Pushing the responsibility for feelings onto the other person)
    When the exhausted husband pushes open the door, if he is met not with
    honesty about feelings, but with a cold accusation: ``You finally know
    how to come back? You don't have this family in your heart at all,
    you're too selfish!''

    These words act like an invisible subpoena, instantly dragging the
    husband to the \textbf{``defendant's seat''} in the relationship court.
    At this moment, the amygdala in the husband's brain sounds the alarm,
    activating a strong \textbf{defense mechanism} \citep{lieberman2007}.
    Sitting in the ``defendant's seat,'' he no longer has the capacity to
    notice his wife's loneliness, because he is busy searching for evidence
    to defend himself:
    ``Who am I working so late for? If I don't work hard overtime, who will
    pay the mortgage? You only see that I come back late, why don't you see
    how much I've given for this family?''

    In this trial of ``who is right and who is wrong,'' both sides try to
    prove that they are the ``victim'' and the other is the ``perpetrator.''
    As the debate intensifies, the wife continues: ``You still dare to
    argue? You don't have the ability to do your job well and you turn
    around and blame me? You don't deserve to be the man of this house!
    Chores, children, I'm doing it all! What have you done?''
    The husband is pushed to his limit and finally explodes: ``You don't
    deserve to be the woman of this house! What kind of family raised a
    shrew like you?''

    If the husband attacks outward, he might physically strike his wife.
    At this moment, conjugal affection turns into mutual hatred.
    The wife truly pushes the husband into the ``defendant's seat,'' and
    the husband goes to prison for domestic violence.
    The wife then ``no longer believes in love,'' raising the child alone.
    And the child, witnessing all this, grows up full of fear and insecurity,
    amidst the mother's curses and resentment toward the husband \citep{vanderkolk2014}.

    If the husband chooses to attack inward, he might choose jackal silence,
    or silently blame himself: ``Why am I so useless? I can't even support a
    family, and I get scolded like this by my wife? I really am a failure!''
    Physically and mentally exhausted, frustrated inside and out, he chooses
    to escape this relationship, leaving this home forever, and eventually
    sinks into illness and bad habits.
    The wife also feels deeper loneliness and despair due to his departure,
    developing psychological illnesses.
    The child, unable to go to school due to lack of money, blames the
    family's incompetence and heartlessness, eventually turning to crime \citep{levine1997}.

    \item \textbf{Empathetic Thinking (Natural Communication):}
    ``I feel lonely, because I really need your companionship.''
    (Stating the need behind the feeling)

    The husband pushes open the door, sees the tears shimmering in his wife's
    eyes, and hears her trembling voice say: ``I feel lonely, because I really
    need your companionship.''
    This sentence is like a key, unlocking each other's hearts.
    The husband holds his wife, and they weep in each other's arms. In this
    embrace, the wife feels a long-lost sense of security and being seen, and
    the husband also feels that he is satisfying his wife's needs bit by bit \citep{porges2011}.

    After calming down, the husband confesses to his wife the stress and
    exhaustion of overtime: ``Recently the company is laying off people, I'm
    under a lot of pressure, afraid of losing our source of income after
    losing my job, so I've been working late. I hope our child can realize
    their dreams, and we can also provide a high-quality life for our family.''
    Hearing her husband's honesty, the wife feels his vulnerability and effort,
    her heart filled with emotion and understanding.
    The two begin to discuss how to jointly meet the need for family financial
    security while also meeting the wife's need for companionship.
    Both sides agree that the husband will come home punctually at 7 PM on
    Mondays and Fridays, and work overtime at other times.

    The husband's company later laid off employees, but because of his
    previous hard work and performance, he successfully transferred to another
    department, continuing to provide stable income for the family.
    Because of her husband's understanding and support, the wife also feels
    seen and valued in this family, her heart filled with security and
    satisfaction.
    The child successfully entered their ideal university and found their
    passion.
    After experiencing this crisis, the family relationship became more solid
    and deep, empathizing and caring for each other even in difficulties.
    Influenced by their parents, the child later formed another family full
    of love, continuing to pass on love and happiness \citep{siegel2003}.
\end{itemize}

Happy families have not avoided misfortune, but when misfortune occurs,
both sides can return to an awareness of their own feelings and needs,
express them honestly, and find solutions that willingly satisfy each other
in the other's needs.
And our way of thinking in the face of misfortune determines different life
trajectories: hate and love are two expressions of the same encounter, and
which mindset we choose determines whether we move toward hate and violence
or toward love and peace \citep{gottman1997}.

\paragraph{The Commonality of Humanity}
Needs are the internal driving forces everyone possesses.
Regardless of race, gender, age, or background, everyone has the same core
needs.
We can imagine an extreme contrast: a CEO in a bustling metropolis and a
shepherd in a remote mountain area.
Their life paths may never cross, and their ways of satisfying needs are
worlds apart---the CEO may gain ``a sense of achievement'' by improving
performance, while the shepherd may gain ``a sense of security'' by
guarding the flock.
However, when we strip away these specific behaviors and touch the deepest
core, you will find that their needs for life are exactly the same; what
differs are the methods of realization \citep{ekman1992}.

This ``commonality'' is the foundation that makes empathy possible.
The reason we can empathize with others' pain is that we have the same
needs in our hearts.
When we can penetrate the other person's angry expressions, accusing words,
and seemingly absurd behavior, and see the universal human needs behind
them that are no different from our own, the other person is no longer an
``unreasonable freak'' or ``enemy,'' but a person just like us who is
trying hard to find satisfaction.
Seeing commonality is the most effective antidote to dissolving ``hostility.''
Focusing on universal human needs is the ``wick'' for us to detach from
dualistic thinking.

\begin{keyinsight}
Our feelings are not caused by the actions of others, but by our own
underlying needs being met or unmet.
By recognizing that all human beings share the same core needs, we can
transform conflict into connection and foster empathy, completely shifting
from judgmental defense to compassionate understanding.
\end{keyinsight}
\section{Strategies}

\subsection{Maslow's Hierarchy of Needs: The Gradient and Dynamic of Needs}

The ``Hierarchy of Needs'' proposed by Abraham Maslow in 1943 is one of
the most influential models in psychology history.
He depicted human needs as a progressive system from basic survival to
advanced spiritual realization:

\begin{itemize}
    \item \textbf{Physiological needs:} Breathing, food, sleep. This is
    the bottom cornerstone of survival \citep{damasio1994}.
    \item \textbf{Safety needs:} Personal security, health guarantee,
    resource ownership \citep{goleman1995}.
    \item \textbf{Love and belonging:} Friendship, love, intimacy.
    This is the core connection most frequently touched upon in communication \citep{ainsworth1978}.
    \item \textbf{Esteem:} Self-respect, sense of achievement, recognition
    by others \citep{erikson1950}.
    \item \textbf{Self-actualization:} Creativity, morality, problem-solving
    ability, unleashing one's potential.
\end{itemize}

\paragraph{Deficit Needs vs. Being Needs}
Maslow collectively referred to the first four levels as \textbf{``Deficit
Needs.''}
They are like the body's ``hunger''; once missing, the individual falls
into tension and anxiety, and all attention is drawn to fill these voids.
And self-actualization is a \textbf{``Being Need,''} it does not originate
from deficiency, but from the desire for life's abundance \citep{salovey1990}.

\paragraph{A Dynamic Flowing Perspective}
Although the pyramid structure seems like steps, Maslow emphasized in his
later years that the satisfaction of needs is \textbf{dynamic and highly
fluid.}
A painter who is willing to endure hunger for artistic creation
(self-actualization) proves that higher-level needs can become the driving
force for lower-level needs \citep{hoffman2000}.
In relationship conflicts, understanding this is crucial:
When the other person shows anger or attack, he is often due to ``esteem''
or ``belonging'' deficit needs being threatened, retreating into a defensive
state.
\textbf{Identifying which level of the pyramid the other person is currently
standing on is the key to providing ``precise empathy.''}

\subsection{Self-Determination Theory}
The Self-Determination Theory (SDT) proposed by Edward Deci and Richard
Ryan argues that three core psychological needs are the basis of human
intrinsic motivation:
\begin{itemize}
    \item \textbf{Autonomy:} Feeling that one's behavior is out of free
    choice, rather than forced obedience \citep{thomas1977}.
    \item \textbf{Competence:} Feeling capable of dealing with challenges
    in the environment and operating effectively \citep{dweck2006}.
    \item \textbf{Relatedness:} Feeling authentically connected with others,
    feeling cared for and caring for others \citep{feldman2024}.
\end{itemize}
Studies have shown that when these three major needs are consistently met
in a relationship, individuals will show higher intrinsic motivation,
better mental health status, and more stable relationship satisfaction.
Conversely, the long-term lack of any one need will lead to psychological
sluggishness, relationship conflict, and inner emptiness.

\begin{table}[h]
\centering
\caption{Universal Human Needs}
\label{tab:universal_needs}
\begin{tabular}{|l|p{11cm}|}
\hline
\textbf{Category} & \textbf{Core Needs} \\ \hline
Connection & Acceptance, belonging, intimacy, love, companionship, communication, understanding, contribution, mutuality, security, celebration \\ \hline
Physical Well-Being & Air, food, water, shelter, touch, rest, movement, sex, safety (physiological), health \\ \hline
Integrity & Authenticity, honesty, consistency \\ \hline
Play & Fun, humor, ease, joy \\ \hline
Peace & Beauty, harmony, inspiration, order, tranquility \\ \hline
Autonomy & Choice, freedom, independence, space, spontaneity \\ \hline
Meaning & Awareness, challenge, clarity, competence, creativity, efficacy, growth, hope, learning, mourning, purpose, self-expression \\ \hline
\end{tabular}
\end{table}

As shown in Table \ref{tab:universal_needs}, human needs are multi-layered
and multi-dimensional.
When we can replace judgments with the identification of needs in daily
communication, many seemingly irreconcilable conflicts will show new
possibilities.

\begin{keyinsight}
Human needs range from basic survival and deficit needs to profound desires
for autonomy, competence, and relatedness.
By recognizing that these fundamental psychological drives shape our internal
motivations and relationship dynamics, we can precisely identify the roots of
conflict and respond with targeted empathy rather than reactive judgment.
\end{keyinsight}
\section{The Nine Faces of Pseudo-Needs}

\subsection{Definition}
A \textbf{strategy} is the specific way, method, or demand we adopt to
satisfy a certain need.
It is specific, personal, and sometimes time-sensitive \citep{gordon1970}.
A \textbf{need} is the deeper, universal internal desire behind the strategy.

For example:
\begin{itemize}
    \item \textbf{Strategy:} ``You must be home by 10 PM every night.''
    \item \textbf{Need:} Security, companionship, connection \citep{kabatzinn1990}.
\end{itemize}

\begin{itemize}
    \item \textbf{Strategy:} ``You have to message me every day saying you
    love me.''
    \item \textbf{Need:} To be valued, to be loved, confirmation \citep{brown2012}.
\end{itemize}

\begin{itemize}
    \item \textbf{Strategy:} ``You cannot eat alone with that colleague.''
    \item \textbf{Need:} Security, trust, loyalty \citep{kohn2005}.
\end{itemize}

When we mistake a strategy for a need, we begin to use various means
(demands, pressure, judgment, punishment) to force the other party to
execute the ``only correct way of satisfaction.''
This often triggers the other party's defense, creating new conflicts,
while the real need remains unmet \citep{faber1980}.

\paragraph{Case Analysis: Confusing Needs and Strategies}

The wife asks her husband to ``come home early every day.''
The husband considers this unreasonable control and feels oppressed and
unfree.
Both sides argue, and the conflict escalates.

But if we peel off the shell of the strategy and explore the needs behind
it:
\begin{itemize}
    \item \textbf{Wife's Needs:} Companionship, connection, security,
    feeling valued in the family \citep{cohen2001}.
    \item \textbf{Husband's Needs:} Autonomy, respect, feeling valued at
    work.
\end{itemize}

These two needs do not conflict!
What conflicts are the strategies---``coming home early'' vs. ``completing
work tasks.''
When both sides see each other's needs clearly, they can jointly find
creative strategies that meet both people's needs simultaneously:
For example, communicating the reasons and times for returning late in
advance, or agreeing on exclusive companionship time \citep{greene1998}.
\textbf{Seeing needs, rather than obsessing over strategies, is the key to
transforming conflict into connection.}

\subsection{Pseudo-Needs}
\begin{table}[h]
\centering
\caption{Common Pseudo-Needs vs. Real Needs}
\label{tab:pseudo_needs}
\begin{tabular}{|p{4.5cm}|p{3.5cm}|p{6.5cm}|}
\hline
\textbf{Pseudo-Need Expression (Strategy/Demand)} & \textbf{Underlying Real Need} & \textbf{Possible Communication Risk} \\ \hline
``I need you to want me every day.'' & Connection, being valued, security & The other person feels forced to execute a task, rather than being emotionally desired. \\ \hline
``I need you to be more responsible.'' & Support, ease, shared burden & Implies personality judgment, highly likely to trigger defense and counterattack. \\ \hline
``I need you to listen to me.'' & Being seen, respect, sense of efficacy & Evolves into a power struggle, ignoring the basis of equal dialogue. \\ \hline
``I need you to give me some space.'' & Autonomy, tranquility, self-regulation & Without background explanation, the other person will feel neglected or abandoned. \\ \hline
``I need you to agree with my views.'' & Being understood, confirmation, harmony & Forcing unified concepts, leading the other person to suppress real thoughts. \\ \hline
``I need you to change your attitude.'' & Respect, security, acceptance & Expression is too vague, the other person feels frustrated due to not knowing what to do. \\ \hline
``I need you to stop being so selfish.'' & Being considered, fairness, mutuality & Pure moral accusation, completely blocking the exploration of needs. \\ \hline
``I need you to share the load.'' & Support, collaboration, fairness & If issued as an order, it weakens the motivation for voluntary contribution. \\ \hline
``I need you to trust me.'' & Connection, acceptance, honesty & Trust is the result of interaction and cannot be obtained through a unilateral demand. \\ \hline
\end{tabular}
\end{table}

As shown in Table \ref{tab:pseudo_needs}, similar to ``pseudo-feelings,''
we also have expressions of ``pseudo-needs''---those words that seem to be
stating needs but are actually expressing judgments or strategic demands.
Identifying pseudo-needs and translating them into expressions of real
needs can fundamentally change the texture of communication---it moves the
focus from ``your behavior'' to ``my heart,'' creating more possibilities
for both sides \citep{siegel2011}.

\paragraph{Psychological Mechanism: Why Do We Habitually Express Pseudo-Needs?}
The reason we often fall into the expression pattern of pseudo-needs usually
originates from deep psychological defense mechanisms and the evasion of
responsibility:
\begin{itemize}
    \item \textbf{Defense Mechanisms and False Sense of Control:} Expressing
    real needs (such as ``I feel insecure and need your support'') makes us
    appear vulnerable, as if handing over the ``power to hurt ourselves''
    to the other party \citep{lazarus1991}.
    In contrast, expressing pseudo-needs (such as ``I need you to be more
    responsible'') is more like issuing an order.
    This condescending posture can bring a false ``sense of control,''
    allowing us to hide behind tough language and avoid facing our inner
    powerlessness.
    \item \textbf{Transferring Responsibility:} Pseudo-needs are essentially
    demanding that the other party be responsible for our feelings.
    This allows us to escape the responsibility of exploring our own needs,
    and instead push the burden of solving the problem entirely onto the
    other party.
    However, this transfer usually triggers the other party's defense
    (such as counterattack or silence), causing the problem to fall into a
    stalemate.
    \item \textbf{The Influence of Culture and Socialization:} From childhood
    to adulthood, we may be educated to be ``good children'' or ``good
    partners,'' instilled with concepts of ``you should...'' and ``you
    must...''.
    This cultural background reinforces our habit of using judgments and
    demands to express needs, rather than directly expressing inner desires.
    \item \textbf{Default Status Differences:} In relationships such as
    parent and child, superior and subordinate, we may be accustomed to
    using imperative language to express needs because it conforms to our
    default status difference in the relationship.
\end{itemize}

When we cannot clearly identify the universal needs deep within our hearts,
these needs often ``deform'' into aggressive or controlling pseudo-needs
(strategies).

\begin{keyinsight}
Confusing our specific strategies for getting what we want with the actual
universal needs that drive us is a primary source of conflict.
When we mistakenly frame demands or behavioral controls as our ``needs,''
we inevitably trigger defensiveness; instead, we must learn to articulate
our true underlying vulnerabilities and desires.
\end{keyinsight}
\section{Skills}

\subsection{Identifying Needs}

Needs are sometimes not easy to perceive directly; they are often hidden
behind feelings, judgments, or strategies \citep{rizzolatti2004}.
Here is a systematic three-step method to help you move from feeling to
need:

\begin{enumerate}
    \item \textbf{Step 1: Feel the present moment.}
    When a strong emotion arises (anger, sadness, fear), do not rush to
    act or express.
    Take a deep breath and bring your attention to your body:
    \textbf{``What am I feeling right now?''} (returning to the skills in
    Chapter 2) \citep{decety2021}.

    \item \textbf{Step 2: Ask from feeling to need.}
    Ask yourself: \textbf{``What unmet need is this feeling telling me
    about?''}
    Refer to the Universal Needs Table (Table \ref{tab:universal_needs})
    and try to name the need that is closest to your heart \citep{warneken2018}.

    \item \textbf{Step 3: Distinguish need from strategy.}
    Check whether the ``need'' you found is truly a need or a strategy.
    Ask yourself: \textbf{``Is this a universal desire that anyone would
    have, or a specific thing I want the other person to do?''}
    If it is the latter, keep digging deeper until you touch the real
    need \citep{streubel2021}.
\end{enumerate}

\begin{table}[h]
\centering
\caption{Correspondence between Feelings and Underlying Needs}
\label{tab:feeling_to_need}
\begin{tabular}{|l|p{5cm}|p{7cm}|}
\hline
\textbf{Feeling} & \textbf{Possible Underlying Need} & \textbf{Example} \\ \hline
Loneliness & Connection, companionship, belonging & I long to feel authentically connected \\ \hline
Anger & Fairness, respect, autonomy & I long for my boundaries to be recognized and respected \\ \hline
Fear & Security, certainty, protection & I long to know that I am safe in this relationship \\ \hline
Sadness & Mourning, companionship & I long for my loss to be authentically seen and accompanied \\ \hline
Anxiety & Order, certainty & I long for a clear understanding of what is going to happen \\ \hline
Grievance & Fairness & I long for my efforts and intentions to be correctly understood \\ \hline
Emptiness & Meaning, connection, sense of value & I long for my existence and actions to be meaningful \\ \hline
Jealousy & Belonging, security, self-worth & I long to be cherished in this relationship \\ \hline
Loss & Contribution seen, connection, expectation & I long for my efforts to be felt and recognized \\ \hline
\end{tabular}
\end{table}

\subsection{``What Do You Need?''}

Similar to ``How are you feeling right now?'' from the previous chapter,
cultivating the habit of asking \textbf{``What do you need right now?''}
from a feeling in interpersonal interactions is one of the most powerful
tools in empathetic communication \citep{lobue2022}.

The power of this question lies in:
\begin{itemize}
    \item It guides the other person's attention from ``whose fault is it''
    to ``what do I truly desire in my heart,'' bypassing the battlefield of
    judgment.
    \item It gives the other person the experience of being valued and seen,
    directly triggering oxytocin secretion at the brain level \citep{christou2025}.
    \item It creates a foundation for both parties to find creative
    solutions---because needs can resonate, while strategies can be
    negotiated \citep{ruba2020}.
\end{itemize}

However, this question is only effective after the other person's emotions
have stabilized (after cortisol levels decrease).
As discussed in Chapter 3, in a state of high cortisol, people cannot
engage in deep introspection---they can only see the other person's
``mistake,'' not anyone's needs \citep{aghaziarati2023}.

\paragraph{When the Other Person Cannot Identify Their Own Needs}
Many people have never been taught to recognize their own needs.
They were trained from childhood to ``obey'' or ``suppress,'' never having
the opportunity to ask themselves ``What do I need?''.
When you ask such a person ``What do you need most?'', they may:
\begin{itemize}
    \item Remain silent or bewildered---indicating that they have never
    been allowed or learned to explore their own needs.
    \item Give a strategy instead of a need---``I need you to come home
    early'' (strategy), rather than ``I need companionship'' (need).
    \item Answer with a judgment---``I need you to stop being so selfish''
    (judgment), rather than ``I need to be valued'' (need).
\end{itemize}

At this time, you can gently assist the other person in translating:
\begin{itemize}
    \item ``When you say you need me to come home early, I think I hear
    that you need more companionship?''
    \item ``When you say you need me to stop being so selfish, I think I
    hear that you need to feel valued and considered?''
\end{itemize}

This \textbf{``Translational Empathy''} is a high-level skill; it is not
to correct the other person, but to help the other person's brain find
deeper language \citep{dunn2000}.

\begin{keyinsight}
Identifying and expressing needs involves systematically moving past
immediate emotional reactions and distinguishing between specific strategies
and fundamental human desires.
When we actively translate judgments or vague requests into the clear language
of needs, we defuse conflict and invite profound collaboration and connection.
\end{keyinsight}
\section{Exercises: Finding Needs}

\subsection{From Feelings to Needs}

\begin{enumerate}
    \item \textbf{``I feel very lonely, as if I am not important.'' --- What is the need behind this?}
    \begin{itemize}
        \item \textbf{Analysis:} Loneliness usually points to unmet needs
        for connection, companionship, and belonging; ``not important''
        points to the need to be seen and valued.
        \item \textbf{Need:} Connection, companionship.
        \item \textbf{Expression Practice:} ``I feel very lonely, because I
        deeply desire to feel that we are authentically connected.''
    \end{itemize}

    \item \textbf{``I get annoyed as soon as I hear her voice.'' --- What is the need behind this?}
    \begin{itemize}
        \item \textbf{Analysis:} Annoyance is often a manifestation of
        elevated cortisol, indicating a core need has been unmet for a long
        time \citep{howard2024}. Here we need to ask: In this relationship,
        what has not been satisfied?
        \item \textbf{Possible Needs:} Respect, autonomy, tranquility,
        connection.
        \item \textbf{Expression Practice:} ``Every time we talk, I feel
        very annoyed, I think it's because I deeply desire to feel equally
        treated in our communication.''
    \end{itemize}

    \item \textbf{``He never thanks me for everything I do.'' --- What is the need behind this?}
    \begin{itemize}
        \item \textbf{Analysis:} ``Never'' is an absolute judgment, hiding
        long-suppressed grievances. The real need is for contribution to
        be seen \citep{izard2001}.
        \item \textbf{Need:} Understanding, fairness.
        \item \textbf{Expression Practice:} ``I feel somewhat lost, because
        I long for my efforts to be understood and treated fairly.''
    \end{itemize}

    \item \textbf{``I am really worried about our future.'' --- What is the need behind this?}
    \begin{itemize}
        \item \textbf{Analysis:} Worry is a signal of rising cortisol,
        warning that a core need involving the future is facing a threat \citep{herman2025}.
        \item \textbf{Need:} Security, clarity, and consensus.
        \item \textbf{Expression Practice:} ``I feel somewhat worried,
        because I deeply desire that we can jointly have a clear direction
        and consensus about our future.''
    \end{itemize}

    \item \textbf{``Every time I give an opinion, he says I'm being nosy.'' --- What is the need behind this?}
    \begin{itemize}
        \item \textbf{Analysis:} Being denied the right to participate
        hides the need for contribution and acceptance \citep{hoffman2000}.
        \item \textbf{Need:} Respect.
        \item \textbf{Expression Practice:} ``Every time I share my
        thoughts, I feel somewhat hurt, because I deeply desire my thoughts
        to be respected.''
    \end{itemize}

    \item \textbf{``I don't know why, but I can't get motivated to do anything.'' --- What is the need behind this?}
    \begin{itemize}
        \item \textbf{Analysis:} This is a typical description of low
        dopamine and a lack of meaning \citep{dweck2006}.
        \item \textbf{Need:} Meaning, purpose, creation, connection, growth.
        \item \textbf{Expression Practice:} ``Recently I feel somewhat dull,
        I think I internally desire to do some things that make me feel
        truly meaningful.''
    \end{itemize}

    \item \textbf{``He always turns a deaf ear to my words.'' --- What is the need behind this?}
    \begin{itemize}
        \item \textbf{Analysis:} ``Always'' is a judgment, hiding the unmet
        needs of being listened to and valued \citep{rogers1961}.
        \item \textbf{Need:} Connection, respect.
        \item \textbf{Expression Practice:} ``Every time I finish speaking
        and get no response, I feel somewhat lost, because I deeply desire
        to feel connected.''
    \end{itemize}

    \item \textbf{``I feel that he is controlling everything about me.'' --- What is the need behind this?}
    \begin{itemize}
        \item \textbf{Analysis:} ``Being controlled'' is a pseudo-feeling,
        hiding a severe lack of the needs for autonomy and freedom \citep{thomas1977}.
        \item \textbf{Need:} Autonomy.
        \item \textbf{Expression Practice:} ``When my schedule and choices
        need to pass your approval, I feel very oppressed, because I deeply
        desire to be able to freely make my own choices.''
    \end{itemize}

    \item \textbf{``Just thinking about going to work tomorrow makes me feel bad.'' --- What is the need behind this?}
    \begin{itemize}
        \item \textbf{Analysis:} This continuous resistance often comes
        from the systemic lack of a core need in the work environment \citep{lazarus1991}.
        \item \textbf{Need:} Meaning, respect, autonomy, belonging, rest \citep{kabatzinn1990}.
        \item \textbf{Expression Practice:} ``Just thinking about tomorrow,
        I feel somewhat heavy. I think what I internally desire is an
        environment that makes me feel my work is meaningful.''
    \end{itemize}
\end{enumerate}

\begin{keyinsight}
Identifying the underlying need behind our own and others' surface feelings
requires intentional practice.
Whether translating anger into a desire for respect or sadness into a plea
for connection, mastering this skill is essential for breaking negative
communication cycles.
\end{keyinsight}
\section{Practice Exercise}
\begin{practiceexercise}
\textbf{Multiple Choice Questions}
\begin{enumerate}
  \item Which of the following statements represents an authentic expression of a need rather than a strategy or pseudo-need \citep{gordon1970}?
    \begin{itemize}
      \item[A)] ``I need you to take out the trash right now.'' \citep{cohen2001}
      \item[B)] ``I need you to be more responsible.'' \citep{kohn2005}
      \item[C)] ``I need to feel supported and have a sense of fairness in our household chores.''
      \item[D)] ``I need you to stop playing video games all evening.''
    \end{itemize}
  \vspace{\fill}
  \item Your teenager comes home angry, slamming the door, and says, ``You always treat me like a baby!'' Which underlying need is most likely driving this outburst \citep{brown2012}?
    \begin{itemize}
      \item[A)] A need for you to stop talking to them. \citep{dweck2006}
      \item[B)] A need for autonomy, trust, and to be respected as a young adult.
      \item[C)] A need for more rules to guide their behavior.
      \item[D)] A need for an apology for ruining their day.
    \end{itemize}
  \vspace{\fill}
  \item When translating a partner's judgment (e.g., ``You're so selfish!'') into a need, what is the best empathetic response \citep{bowlby1969}?
    \begin{itemize}
      \item[A)] ``I am not selfish; I do plenty for this family.''
      \item[B)] ``Are you feeling upset because you need to feel considered and valued?''
      \item[C)] ``Well, you're selfish too sometimes.''
      \item[D)] ``You need to stop criticizing me.''
    \end{itemize}
\end{enumerate}
\vspace{\fill}
\newpage

\textbf{Explanations for Questions 1--3}

\textbf{Question 1: C.} This statement clearly articulates universal human desires (support, fairness) without attacking the other person's character or demanding a specific, controlling action \citep{dunn2000}. Options A and D are strategies disguised as needs \citep{kabatzinn1990}, and Option B is a judgment disguised as a need \citep{rogers1961}.

\textbf{Question 2: B.} Anger and accusations of being treated ``like a baby'' typically stem from a teenager's thwarted desire for independence and autonomy \citep{thomas1977}. Recognizing this allows the parent to bypass the defensive tone and address the core issue of trust and growing independence \citep{lazarus1991}. A failure to address this can limit proper emotional coaching \citep{gottman1997}.

\textbf{Question 3: B.} This response perfectly demonstrates empathetic translation \citep{siegel2011}. It bypasses the judgmental attack (``selfish'') and checks in with the partner's underlying feeling (upset) and core need (to be considered/valued) \citep{rosenberg2003}, inviting connection rather than starting an argument \citep{greene1998}.
\vspace{\fill}
\newpage
\end{practiceexercise}
\section{Reflection Questions}
\begin{reflection}
\small
\begin{enumerate}[itemsep=0pt, parsep=0pt, topsep=0pt]
  \item When you experience intense negative emotions, do you tend to look inward at your own unmet needs \citep{rizzolatti2004}, or outward to blame someone else?
  \item Which ``pseudo-needs'' (e.g., ``I need you to listen,'' ``I need you to change'') do you catch yourself using most often \citep{streubel2021}?
  \item Think of a recent conflict with your child. What specific strategy were you fighting over, and what was your actual underlying need \citep{lobue2022}?
  \item What fears arise when you consider expressing your deep, vulnerable needs (like a need for security) rather than using demands or judgments \citep{christou2025}?
  \item How did your own caregivers handle their unmet needs? Did they express them clearly, or did they rely on guilt and pseudo-needs \citep{ruba2020}?
  \item In what situations is it hardest for you to separate the behavior (the strategy) from the universal human need driving it \citep{aghaziarati2023}?
  \item How might your relationship change if you responded to your child's anger by asking yourself, ``What are they needing right now?'' \citep{izard2001}
  \item What is one small step you can take today to ensure that your own ``deficit needs'' (like sleep or basic support) are adequately met \citep{herman2025}?
  \item How does practicing self-empathy regarding your own unmet needs impact your ability to be patient with your child's unmet needs \citep{hoffman2000}?
\end{enumerate}
\end{reflection}
